\subsection{Důležité parametry a jejich vliv na herní pocit}

Každý parametr v \texttt{PlayerControllerData} má specifický vliv na chování hráče a celkový herní pocit. Pochopení jejich vzájemného působení je klíčové pro správné nastavení mechanik.

Mezi nejdůležitější parametry běhu patří \texttt{runMaxSpeed}, který definuje maximální rychlost horizontálního pohybu. Parametr \texttt{runAcceleration} určuje, jak rychle hráč dosáhne této maximální rychlosti. Vysoká hodnota zrychlení vede k okamžitému startu, zatímco nízká hodnota vytváří efekt rozjezdu. Důležitým parametrem je také \texttt{accelInAir}, který určuje procentuální poměr pozemního zrychlení použitého ve vzduchu. Nízká hodnota tohoto parametru ztěžuje změnu směru během skoku, což může být žádoucí pro obtížnější sekce.

Pro skok jsou klíčové parametry \texttt{jumpForce} určující sílu počátečního výskoku, \texttt{coyoteTime} definující dobu grace period po opuštění platformy, a \texttt{jumpBufferTime} pro registraci skoku před dopadem. Parametr \texttt{jumpHangTimeThreshold} určuje prahovou rychlost, při které se aktivuje efekt vznášení na vrcholu skoku.

Gravitace je řízena parametrem \texttt{gravity} pro základní gravitační sílu a \texttt{fallMultiplier} pro zrychlení pádu. Vyšší hodnota fall multiplieru vede k rychlejšímu pádu a kratšímu času ve vzduchu, což zlepšuje kontrolu nad přesným načasováním dopadu.

Dash je konfigurován pomocí \texttt{dashSpeed} pro rychlost výpadu, \texttt{dashTime} pro dobu trvání a \texttt{dashCooldown} pro prodlevu mezi dashi.